\chapter{От плоской волны к конечномерной задаче}
\label{math:fem}
\label{math:chapter}

\section{Падающее, рассеянное и полное поля}
Пусть плоская волна распространяется в однородной среде и встречает
включение с другими свойствами. Мы знаем амплитуду и фазу падающей
волны в каждой точке. Требуется найти, как включение изменяет поле.
Для точки $x=(x_1,x_2)$, волнового числа $k>0$ и угла $\theta$
положим
\begin{align}
 d&=(\cos\theta,\sin\theta),\quad |d|=1,\label{math:direction}\\
 u_i(x)&=\exp(\ii k d\cdot x)
       =\cos(kd\cdot x)+\ii\sin(kd\cdot x).\label{math:incident}
\end{align}
Вектор $d$ задаёт направление распространения. Линии равной фазы
перпендикулярны $d$ и повторяются через длину волны $2\pi/k$ вдоль
этого направления. При принятой временно́й зависимости физическое
колебание имеет вид $\operatorname{Re}(u(x)e^{-\ii\omega t})$.
Модуль комплексного $u$ даёт амплитуду, а аргумент при $u\ne0$ даёт фазу.

Полное поле $u$ складывается из известного падающего поля $u_i$ и
искомого рассеянного поля $u_s$. Именно $u_s$ приближает базовая
нейросетевая модель. Зная приближение $u_{s,h}$, можно вычислить полное
поле в любой точке элемента как $u_i(x)+u_{s,h}(x)$. При интегрировании
этой суммы аналитическую экспоненту вычисляют в точках квадратуры.
Линейная интерполяция суммы узловых значений дала бы другую функцию.

\section{Уравнение для рассеянной волны}
Вычислительная область $\Omega=(-2,2)^2$ окружает неоднородность.
Положительные скалярные коэффициенты $A,n$ ограничены сверху и отделены
от нуля, а в фоновой среде равны единице. Полное поле удовлетворяет
\begin{equation}
 \nabla\cdot(A\nabla u)+k^2nu=0.\label{math:total-pde}
\end{equation}
Для падающего поля $\nabla u_i=\ii k d\,u_i$, откуда
$\Delta u_i+k^2u_i=0$. Подставим $u=u_i+u_s$ в уравнение и воспользуемся
этим тождеством. Тогда
\begin{equation}
 \nabla\cdot(A\nabla u_s)+k^2nu_s
 =\nabla\cdot((1-A)\nabla u_i)+k^2(1-n)u_i.\label{math:scattered-pde}
\end{equation}
Контрасты $1-A$ и $1-n$ задают источник рассеянной волны.
При скачках коэффициентов дивергенцию понимают в смысле распределений.

Чтобы ограничиться квадратом, на его границе $\Gamma=\partial\Omega$
задают приближённое условие выхода волны первого порядка
\begin{equation}
 \partial_\nu u_s=\ii k u_s,\qquad
 \partial_\nu u_s=\nabla u_s\cdot\nu,\label{math:robin}
\end{equation}
где $\nu$ есть внешняя единичная нормаль на стороне квадрата. На данной
стороне это условие точно выполняется для плоской волны, выходящей по её
нормали.
При косом выходе нормальная составляющая волнового вектора меньше $k$,
поэтому искусственная граница может отражать часть волны. Это источник
погрешности ещё до построения сетки.

\section{Слабая форма}
Для вывода обозначим $q=(1-A)\nabla u_i$. Умножим
\eqref{math:scattered-pde} на $\bar v$, где $v\in H^1(\Omega;\C)$,
и проинтегрируем по частям. Черта означает комплексное сопряжение.
Для гладких полей формула Грина даёт
\begin{align}
 -\int_\Omega A\nabla u_s\cdot\nabla\bar v\,dx
 +\int_\Gamma A\partial_\nu u_s\bar v\,ds
 +k^2\int_\Omega nu_s\bar v\,dx
 ={}&-\int_\Omega q\cdot\nabla\bar v\,dx\notag\\
 &+\int_\Gamma(q\cdot\nu)\bar v\,ds
 +k^2\int_\Omega(1-n)u_i\bar v\,dx.\label{math:green}
\end{align}
В окрестности границы $A=1$, поэтому граничный вклад от $q$ исчезает.
Подставив \eqref{math:robin} и умножив обе части на $-1$, получаем
слабую постановку. Найти $u_s\in H^1(\Omega;\C)$, для которого при
каждом $v\in H^1(\Omega;\C)$ выполнено
\begin{align}
 a(u_s,v)&=\ell(v),\label{math:weak}\\
 a(w,v)&=\int_\Omega A\nabla w\cdot\nabla\bar v\,dx
 -k^2\int_\Omega nw\bar v\,dx
 -\ii k\int_\Gamma w\bar v\,ds,\label{math:bilinear}\\
 \ell(v)&=\int_\Omega(1-A)\nabla u_i\cdot\nabla\bar v\,dx
 +k^2\int_\Omega(n-1)u_i\bar v\,dx.\label{math:rhs}
\end{align}
Форма линейна по $w$ и сопряжённо-линейна по $v$. В интегралах
содержатся только первые производные поля и тестовой функции. Такая
запись определена и для кусочно постоянного $A$ со скачками на границах
включений.

\begin{worked}{Проверка правой части}
В однородной среде $A=n=1$, поэтому $\ell(v)=0$ для любого $v$.
При единственности решения имеем $u_s=0$ во всей области, а полное поле
совпадает с падающим. Эта проверка обнаруживает паразитный источник.
Для проверки знака массового вклада возьмём $A=1$ и $n\ge1$, причём
$n>1$ на множестве положительной площади. При $v=u_i$ получаем
$\ell(u_i)=k^2\int_\Omega(n-1)\,dx>0$, поскольку $|u_i|^2=1$.
\end{worked}

\section{Метод конечных элементов на треугольниках}
Разобьём область на треугольники. Узловая базисная функция $\varphi_j$
линейна на каждом элементе, непрерывна и удовлетворяет
$\varphi_j(x_i)=\delta_{ij}$. Пусть $N$ есть число узлов. Ищем приближение
\begin{equation}
 u_{s,h}(x)=\sum_{j=1}^N U_j\varphi_j(x),\qquad
 U=(U_1,\ldots,U_N)^\top\in\C^N.\label{math:fem-expansion}
\end{equation}
Подставим это разложение в \eqref{math:weak} и последовательно возьмём
$v=\varphi_i$, $i=1,\ldots,N$. Получим систему
\begin{align}
 (K-k^2M_n-\ii k B_\Gamma)U&=F,\label{math:fem-system}\\
 K_{ij}&=\int_\Omega A\nabla\varphi_j\cdot\nabla\varphi_i\,dx,
 & (M_n)_{ij}&=\int_\Omega n\varphi_j\varphi_i\,dx,\notag\\
 (B_\Gamma)_{ij}&=\int_\Gamma\varphi_j\varphi_i\,ds,
 & F_i&=\ell(\varphi_i).\notag
\end{align}
Индекс $i$ нумерует уравнения, $j$ нумерует неизвестные.
Матрицы $K,M_n,B_\Gamma$ вещественные и симметричные, каждая имеет
размер $N\times N$. Их комбинация в \eqref{math:fem-system}
комплексно-симметрична и, при ненулевой граничной матрице, неэрмитова.
При неизменных $A,n,k$ и сетке смена направления $d$ меняет только
нагрузку $F\in\C^N$. Для нескольких направлений достаточно одной
факторизации матрицы.

\subsection{Вклад одного элемента}
Пусть на треугольнике $T$ площади $|T|$ коэффициенты равны $A_T,n_T$.
Локальные матрицы жёсткости и массы получают умножением на $A_T$ и $n_T$
соответственно следующих матриц:
\begin{equation}
 (K_T^{(0)})_{ij}=|T|\nabla\varphi_i\cdot\nabla\varphi_j,\qquad
 M_T^{(0)}=\frac{|T|}{12}
 \begin{pmatrix}2&1&1\\1&2&1\\1&1&2\end{pmatrix}.\label{math:local-mass}
\end{equation}
На граничном ребре $e$ длины $L_e$ интегрирование даёт
\begin{equation}
 B_e=\frac{L_e}{6}\begin{pmatrix}2&1\\1&2\end{pmatrix}.\label{math:edge-mass}
\end{equation}
Каждый локальный элемент добавляют к строкам и столбцам общих узлов.
Для правой части из \eqref{math:rhs} аналогично следует
\begin{equation}
 F_{T,j}=(1-A_T)\ii k(d\cdot\nabla\varphi_j)\int_T u_i\,dx
              + k^2(n_T-1)\int_Tu_i\varphi_j\,dx.\label{math:local-rhs}
\end{equation}
В отличие от произведений линейных функций, здесь под знаком интеграла
стоит экспонента. В используемом решателе эти интегралы вычисляются
приближённо, по квадратуре Гаусса с $8\times8$ точками после
преобразования Даффи из квадрата в треугольник.

\begin{worked}{Проверка на единичном прямоугольном треугольнике}
Для вершин $(0,0),(1,0),(0,1)$ базисные функции равны $1-x_1-x_2$, $x_1$,
$x_2$, а площадь равна $1/2$. Их градиенты дают
\[
 K_T^{(0)}=\tfrac12\begin{pmatrix}2&-1&-1\\-1&1&0\\-1&0&1\end{pmatrix}.
\]
Сумма элементов каждой строки равна нулю. Поэтому постоянное поле с
вектором значений $\bm1=(1,1,1)^\top$ имеет нулевую энергию градиента.
Для матрицы массы проверка другая:
$\bm1^\top M_T^{(0)}\bm1=12/24=1/2=\int_T1\,dx$.
\end{worked}

\section{Как измерить ошибку поля}
Пусть $M$ собрана из локальных матриц $M_T^{(0)}$ без коэффициента $n_T$.
Это согласованная матрица массы. Для комплексного вектора $z$ и
соответствующей ему конечно-элементной функции $z_h$ имеем
\begin{equation}
 \|z_h\|_{L^2(\Omega)}^2=z^*Mz,\label{math:consistent-norm}
\end{equation}
где звёздочка обозначает сопряжённое транспонирование. При построении
базиса удобно использовать диагональную матрицу весов $W$, полученную
суммированием строк $M$,
\begin{equation}
 w=M\bm1,\quad W=\diag(w_1,\ldots,w_N),\qquad
 \|z\|_W^2=z^*Wz=\sum_jw_j|z_j|^2.\label{math:lumped-norm}
\end{equation}
Каждый треугольник добавляет $|T|/3$ в вес каждой вершины. Поэтому на
невырожденной сетке $w_j>0$ и $\sum_jw_j=|\Omega|$. Узлы в сгущённой
части сетки получают меньший вес, чем узлы с большой прилегающей площадью.

На предыдущем треугольнике $W_T=I_3/6$. При $z=(1,0,0)^\top$ квадрат
точной нормы равен $1/12$, а квадрат диагональной нормы равен $1/6$.
Обе формы верно интегрируют квадрат константы, однако между узлами
квадрат линейной функции меняется квадратично.

Для $z=a+\ii b$, где $a,b\in\R^N$, симметрия вещественной $M$ даёт
\begin{equation}
 z^*Mz=a^\top Ma+b^\top Mb.\label{math:complex-norm}
\end{equation}
Следовательно, относительная ошибка прогноза $\widehat U$ для
ненулевого эталона $U$ равна
\begin{equation}
 e_M(U,\widehat U)=
 \frac{\|\widehat U-U\|_M}{\|U\|_M}.\label{math:relative-error}
\end{equation}
В этой дроби обе части поля учитываются до извлечения квадратного корня.
При $U=0$ относительная ошибка математически не определена. Для отчётов
принято возвращать ноль при нулевом прогнозе и бесконечность при ненулевом.
Абсолютная ошибка $\|\widehat U-U\|_M$ определена в обоих случаях.

\begin{idea}{Какие данные описывают один расчёт}
Эталон получают, собирая систему по \emph{ячеечным} значениям $A_T,n_T$.
На вход базовой сети поступают \emph{узловые} значения материалов,
полученные арифметическим усреднением по соседним ячейкам. На границе
включения такое усреднение сглаживает скачок. Одно обучающее наблюдение
связывает эти входы и параметры волны с эталонным узловым вектором $U$.
\end{idea}

\chapter{Сжатие поля с помощью POD и SVD}
\label{math:pod}

\section{Среднее поле и пространственные моды}
Пусть $f^{(1)},\ldots,f^{(m)}\in\R^N$ содержат узловые значения
одной вещественной величины на общей сетке. Каждый вектор называют
\emph{снимком}. Например, можно собрать снимки вещественной части
рассеянного поля при разных углах падения. Требуется найти $r$ общих
пространственных функций и представить каждый снимок в виде
\begin{equation}
 f^{(s)}\approx\mu+\Phi c^{(s)},\qquad
 \mu=\frac1m\sum_{s=1}^m f^{(s)},\qquad
 \Phi=(\phi_1,\ldots,\phi_r)\in\R^{N\times r}.\label{math:pod-expansion}
\end{equation}
Здесь $\mu\in\R^N$ хранит среднее по снимкам. Каждый столбец $\phi_j$
задаёт значения одной моды на всей сетке, а $j$-я координата
$c^{(s)}\in\R^r$ задаёт амплитуду этой моды в снимке $s$.
При $r\ll N$ для описания нового поля в том же пространстве достаточно
$r$ чисел. Построение оптимальных по совокупности снимков мод называется
методом главных ортогональных мод (POD, \emph{proper orthogonal
decomposition}). Сингулярное разложение матрицы (SVD, \emph{singular
value decomposition}) даёт способ их вычислить.

\section{От матрицы снимков к базису}
В базовой модели вещественные снимки расположены по строкам матрицы $F$.
Поэтому центрирование и взвешивание имеют вид
\begin{align}
 F&=\begin{pmatrix}(f^{(1)})^\top\\[-2pt]\vdots\\(f^{(m)})^\top\end{pmatrix}
       \in\R^{m\times N},\qquad
 X=F-\bm1_m\mu^\top,\label{math:row-snapshots}\\
 Y&=XW^{1/2},\qquad
 W^{1/2}=\diag(\sqrt{w_1},\ldots,\sqrt{w_N}).\label{math:weighted-snapshots}
\end{align}
В каждом столбце $X$ сумма равна нулю. Поэтому ранг центрированных данных
не превосходит $\min(m-1,N)$. Положительные веса из
\eqref{math:lumped-norm} учитывают площадь, приходящуюся на узел.
Квадратный корень нужен для следующего равенства:
\begin{equation}
 \sum_{s=1}^m\|f^{(s)}-\widehat f^{(s)}\|_W^2
 =\|(F-\widehat F)W^{1/2}\|_F^2.\label{math:weighted-objective}
\end{equation}
Квадрат нормы Фробениуса $\|Y\|_F^2$ есть сумма квадратов модулей всех
элементов матрицы. Взвешивание стоит справа, поскольку номер узла
нумерует столбцы.

Возьмём экономное сингулярное разложение с $q=\min(m,N)$:
\begin{equation}
 Y=U\Sigma V^\top,\qquad
 U\in\R^{m\times q},\quad\Sigma\in\R^{q\times q},\quad
 V\in\R^{N\times q}.\label{math:svd}
\end{equation}
Здесь буква $U$ обозначает матрицу левых сингулярных векторов. Обе
матрицы сингулярных векторов имеют ортонормированные столбцы.
Числа $\sigma_1\ge\ldots\ge\sigma_q\ge0$ на диагонали $\Sigma$
упорядочены по убыванию. Выберем первые $r$ правых сингулярных векторов
$V_r$ и вернёмся к исходным координатам узлов:
\begin{align}
 \Phi&=W^{-1/2}V_r,\label{math:modes}\\
 \Phi^\top W\Phi&=V_r^\top V_r=I_r.\label{math:orthonormal}
\end{align}
Таким образом, столбцы $\Phi$ ортонормированы в метрике $W$.

\subsection{Коэффициенты известного поля}
У известного поля $f$ ищем ближайшее приближение вида $\mu+\Phi c$.
Остаток должен быть $W$-ортогонален каждому столбцу $\Phi$. Из
$\Phi^\top W(f-\mu-\Phi c)=0$ и \eqref{math:orthonormal} следует
\begin{align}
 c_W&=\Phi^\top W(f-\mu),\label{math:encode-column}\\
 P_W f&=\mu+\Phi c_W.\label{math:decode-column}
\end{align}
Вся таблица коэффициентов и восстановленных снимков равна
\begin{align}
 C&=XW\Phi\in\R^{m\times r},\label{math:encode-row}\\
 F_r&=\bm1_m\mu^\top+C\Phi^\top\in\R^{m\times N}.\label{math:decode-row}
\end{align}
На обучающих снимках $C=YV_r=U_r\Sigma_r$. Для нового, ещё неизвестного
рассеянного поля формула \eqref{math:encode-column} недоступна.
Сеть получает параметры задачи и предсказывает коэффициенты, после чего
восстанавливается узловой вектор.

\section{Пример с двумя узлами}
Возьмём три снимка с двумя координатами и назначим второму узлу
вчетверо больший вес:
\begin{equation}
 F=\begin{pmatrix}0&0\\1&1\\2&2\end{pmatrix},\qquad W=\diag(1,4).
 \label{math:example-data}
\end{equation}
Вычитая среднее $\mu=(1,1)^\top$ и умножая вторую координату на $2$,
получаем
\[
 Y=\begin{pmatrix}-1&-2\\0&0\\1&2\end{pmatrix}
 =\begin{pmatrix}-1\\0\\1\end{pmatrix}\begin{pmatrix}1&2\end{pmatrix}.
\]
Матрица имеет ранг один. Единственное положительное сингулярное число
равно произведению длин этих двух векторов, $\sigma_1=\sqrt{10}$.
Правый сингулярный вектор равен $(1,2)^\top/\sqrt5$. Поэтому мода
$\phi=W^{-1/2}V_1=(1,1)^\top/\sqrt5$ и
$\phi^\top W\phi=(1+4)/5=1$. Три коэффициента равны $-\sqrt5$, $0$ и
$\sqrt5$. Например, из последнего получаем
$\mu+\sqrt5\,\phi=(2,2)^\top$, как в последней строке $F$.

\begin{worked}{Задача. Спроецируйте новое поле}
Для $f=(2,0)^\top$ найдите ближайшее поле $P_W f$ на той же прямой.
Проверьте ортогональность остатка.

\textbf{Решение.} Отклонение от среднего равно $(1,-1)^\top$, поэтому
$c_W=(1-4)/\sqrt5=-3/\sqrt5$. Реконструкция имеет равные координаты
$P_W f=(2/5,2/5)^\top$. Остаток $e=(8/5,-2/5)^\top$ удовлетворяет
$\phi^\top We=(8/5-8/5)/\sqrt5=0$. Квадрат его нормы равен $16/5$.
Второй узел получил больший вес, и общее значение $2/5$ оказалось
ближе к нулю во втором узле, чем к двум в первом.
\end{worked}

\section{Что оптимизирует усечение SVD}
Для точного SVD теорема Эккарта--Янга даёт
\begin{equation}
 \min_{\operatorname{rank}(Z)\le r}\|Y-Z\|_F^2
 =\sum_{j>r}\sigma_j^2.\label{math:eckart-young}
\end{equation}
В силу \eqref{math:weighted-objective} это минимальная сумма квадратов
\emph{абсолютных} ошибок на обучающих снимках в метрике $W$ при фиксированном
среднем. Для $E=\|Y\|_F^2>0$ доля сохранённой энергии отклонений равна
\begin{equation}
 \eta_r=\frac{\sum_{j\le r}\sigma_j^2}{E},\qquad
 \frac{\|(F-F_r)W^{1/2}\|_F}{\|XW^{1/2}\|_F}
 =\sqrt{1-\eta_r}.\label{math:energy}
\end{equation}
Например, $99\,\%$ энергии дают $10\,\%$ относительной ошибки для всей
матрицы \emph{центрированных} снимков в этой норме. У отдельного поля
со своим знаменателем ошибка может оказаться больше или меньше.
При $E=0$ все снимки совпадают со средним и достаточно ранга ноль.

Для произвольных предсказанных коэффициентов $\widehat c$ ортогональность
остатка даёт точное разложение
\begin{equation}
 \|f-\mu-\Phi\widehat c\|_W^2
 =\|f-P_W f\|_W^2+\|c_W-\widehat c\|_2^2.\label{math:pythagoras}
\end{equation}
Первый член уже фиксирован выбранным пространством. Обучение сети может
изменить второй. Если качество поля оценивают в точной конечно-элементной
норме $M$, ближайшие коэффициенты определяют заново:
\begin{equation}
 Gc_M=\Phi^*M(f-\mu),\qquad G=\Phi^*M\Phi.\label{math:consistent-projection}
\end{equation}
При независимых столбцах $\Phi$ и положительно определённой $M$ матрица
Грама $G$ положительно определена. Для $P_M f=\mu+\Phi c_M$ из той же
ортогональности следует разложение с $\|c_M-\widehat c\|_G^2$ вместо
евклидовой нормы коэффициентов. Именно $\|f-P_M f\|_M$ даёт наименьшую
достижимую ошибку поля в выбранном пространстве. Сравнение обеих
проекций доказано в разделе~\ref{rig:projection}.

\section{Два способа представить комплексное поле}
В базовой модели четыре вещественных базиса строятся отдельно для входных
материалов $A,n$ и для выходных полей $\operatorname{Re}u_s$,
$\operatorname{Im}u_s$. У каждого свое среднее, свой ранг и своя таблица
мод. Совпадение номера столбца между двумя таблицами само по себе
не связывает пространственные функции. При прогнозе поле собирается как
\begin{equation}
 \widehat u_s=\mu_R+\Phi_R\widehat c_R+
             \ii(\mu_I+\Phi_I\widehat c_I).\label{math:complex-decode}
\end{equation}
При этом $\widehat c_R\in\R^{r_R}$, $\widehat c_I\in\R^{r_I}$,
а ошибка поля учитывает обе компоненты по \eqref{math:complex-norm}.

Для комплексного POD удобно расположить \emph{центрированные снимки по
столбцам} матрицы $S\in\C^{N\times m}$. Пусть $M=R^*R$ положительно
определена, например $R$ есть верхний треугольный множитель Холецкого.
Если $RS=Q\Sigma V^*$, то
\[
 \Phi=R^{-1}Q_r,\qquad \Phi^*M\Phi=I_r,\qquad
 c=\Phi^*M(f-\mu)\in\C^r.
\]
Теперь одна мода может одновременно описывать изменения амплитуды и фазы.
Знак вещественной моды произволен. Для комплексной моды можно изменить
фазу $\phi_j\mapsto e^{\ii\alpha}\phi_j$ и компенсировать её заменой
$c_j\mapsto e^{-\ii\alpha}c_j$. При кратном сингулярном числе допустим
унитарный поворот внутри соответствующего подпространства. Поэтому при
сравнении вычислений проверяют поля или проекторы.

\section{Как выбрать ранг по данным}
Материалы одной среды могут повторяться при разных углах и частотах. В
обучающую матрицу материалов берут по одному снимку от каждой среды,
а для выходных полей используют все обучающие решения. При фиксированном
материале его центрированные снимки нулевые, и входной ранг равен нулю.
Тогда известную среду представляет одно среднее. Информацию о меняющихся
частоте и направлении передают оставшиеся входные координаты.

Для выходных полей базовая модель перебирает ранги от единицы до заданного
предела. Для обучения нужна хотя бы одна выходная мода. Если обе
компоненты постоянны по снимкам, каждую из них уже описывает её среднее.
При каждом ранге известные поля проверочной выборки проецируют в метрике
$W$, восстанавливают обе компоненты и усредняют относительную ошибку
\eqref{math:relative-error} в метрике $M$. Берут первый ранг, при котором
средняя ошибка не превышает заданного порога. Если такого ранга нет,
оставляют наибольший доступный и отмечают, что порог не достигнут.
У компоненты с меньшим численным рангом сохраняют только имеющиеся моды.
Нулевые столбцы, которые могут добавлять при хранении, в матрицу Грама
\eqref{math:consistent-projection} не включают.

\begin{idea}{Три числа для трёх вопросов}
Остаток по спектру \eqref{math:energy} показывает, сколько вариации потеряно
при сжатии обучающих данных. Ошибка проекции проверочного поля показывает,
насколько подходит пространство для новых условий. Ошибка прогноза на
независимом тесте дополнительно учитывает выбор коэффициентов нейросетью.
Среднее и моды следует вычислять только по обучающим данным. Включение
тестового поля в SVD заранее сообщает методу о пространственных
особенностях этого поля.
\end{idea}

\subsection{Численный ранг и приближённое разложение}
Полное экономное SVD вычисляет все $q$ сингулярных чисел. Для большого
числа снимков можно приближённо найти ведущие направления. Пусть $t$
есть запрошенный ранг, ограниченный сверху числом $q$. При $t=0$ или
нулевой энергии центрированных снимков базис сразу оставляют пустым.
В остальных случаях при $t+10<q$
используемый случайный алгоритм берёт гауссову матрицу
$\Omega_g\in\R^{N\times(t+10)}$ и строит ортонормированный базис $Q$
по столбцам $Y\Omega_g$. Чтобы лучше выделить ведущие направления,
дважды выполняются шаги
\begin{equation}
 Z=\operatorname{orth}(Y^\top Q),\qquad
 Q=\operatorname{orth}(YZ),\label{math:power}
\end{equation}
где $\operatorname{orth}$ оставляет ортонормированные столбцы из
экономного QR-разложения. Затем находят SVD меньшей матрицы $Q^\top Y$.
При $t+10\ge q$ применяется полное экономное разложение.
У приближённых мод измеряют фактическую ошибку восстановления:
равенство с оптимальным хвостом спектра доказано для точного SVD.

Для обоих алгоритмов отсеивают сингулярные числа ниже численного
разрешения. В данной реализации сохраняют числа, которые строго
превышают
\begin{equation}
 \tau=\epsilon_{64}\max(m,N)\max\bigl(\sigma_1,\|FW^{1/2}\|_F\bigr).
 \label{math:numeric-threshold}
\end{equation}
Здесь $\epsilon_{64}$ есть машинная точность, а $\sigma_1$ берётся из
вычисленного разложения. Масштаб до вычитания среднего важен для почти
постоянных полей: малое отклонение от большого фона может быть целиком
обусловлено округлением.
